\section{INTRODUCTION}
\label{ch:introduction}

\iffalse

模糊測試工具（fuzzer）的核心方法是向目標程式（target program）快速提供大量隨機的輸入，這些輸入對於目標程式可能是格式錯誤或非預期的，用以發現目標程式中的漏洞、錯誤或異常行為。透過模糊測試，可以檢測出潛在的安全風險，例如緩衝區溢位、程式崩潰等，這些問題都可能成為攻擊者利用的切入點。近年來，隨著研究不斷地發展，模糊測試的有效性已在多個案例中得到證實，其中尤以 AFL（American Fuzzy Lop）\cite{zalewski2017afl} 最為重要且廣為人知，許多最先進的模糊測試工具皆基於 AFL 修改、實作 \cite{fioraldi2020afl++,pham2019smart,bohme2017directed}。

\if false
模糊測試的本質是一種基因算法（Genetic Algorithms, GAs）\cite{holland1992genetic}，是一類靈感源自自然進化過程的最佳化隨機搜索技術，擅長在廣闊的輸入空間（input space）中最佳化目標問題的答案。有許多研究指出，測試覆蓋率（coverage）為軟體測試的一個重要的指標，越高的覆蓋率通常也代表有更高的機率找出問題 \cite{kochhar2015code,bohme2022reliability}。包含 AFL 在內的一類模糊測試工具，將模糊測試視為找到一組輸入，使測試覆蓋率最大化的最佳化問題。此類方法也被稱為覆蓋率導向模糊測試（Coverage-Guided Fuzzing）。

模糊測試的過程中，每一個提供給目標程式的輸入也稱為種子（seed）。模糊測試利用基因演算法多次的繁衍與篩選種子，生成多樣化的測試輸入。基因演算法根據適應值函數（fitness function）評估一個種子的價值，優勝劣敗。對於覆蓋率導向模糊測試而言，一個種子的適應值是其能夠提供目標程式的覆蓋率。隨著測試的進行，種子池（seed pool）中的種子會越來越多，也代表測試覆蓋率的提升。每當有新的種子生成就會進行適應值評估，若無法提供更多的覆蓋率，便會將其淘汰，反之加入種子池中；適者生存，不適者淘汰。

相較於單純隨機搜索，基因演算法利用「交叉（crossover）、變異（mutation）」等操作產生新的候選解。在 AFL 中，也採用了交叉、變異的方式繁衍種子個體，所有產生的種子皆是初始種子（initial seed）的後代。一個種子代表著一種程式的執行路徑，透過變動父代種子產生子代，期望新的種子能夠因此有不同的執行路徑，進而產生新的程式覆蓋率，這就是覆蓋率導向模糊測試的原理。
\fi

一般的模糊測試工具只專注在輸入檔（input file）和標準輸入（standard input）的變異。當目標程式為命令列工具（command-line tool）時，只能使用固定的命令列選項（command-line option）進行測試。有研究顯示，同時使用多種命令列選項進行模糊測試，可以顯著地增加測試覆蓋率，稱為「多參數模糊測試」（Multi-Argument Fuzzing）或「多選項模糊測試」（Multiple Options Fuzzing） \cite{wu2020sq,tseng2023wei,lu2022guan,tsai2024Multiple} 。此類測試工具提供一個選項組態檔以描述目標程式可以使用的選項，並將選項隨機組合進行測試。

目前的「多參數模糊測試」工具主要有兩個問題。其一是需專家介入，協助撰寫選項組態檔，這是費時費力而且容易出錯的。其次是組態檔仍然只能描述固定、有限的選項，並非完整的選項輸入空間，這可能會導致某些執行路徑難以被探索。雖然已有研究使用大型語言模型（Large Language Model, LLM）協助產生參數組態檔 \cite{tsai2024Multiple}，然而微調（fine-tune）大型語言模型仍然需要人力與時間；大型語言模型透過閱讀目標程式的使用手冊，可以產出對應的組態檔，卻還是缺乏完整探索輸入空間的能力。本研究提出一種新的多參數模糊測試架構，以自動化的方法探索參數的輸入空間，並且提供開源的原型（prototype）。

\else
The core method of a fuzzing tool is to rapidly provide a target program with a large amount of random input, which may be malformed or unexpected, in order to discover vulnerabilities, errors, or abnormal behavior. Through fuzzing, potential security risks can be detected, such as buffer overflows and program crashes, which can become entry points for attackers. In recent years, with continuous research development, the effectiveness of fuzzing has been proven in numerous cases, among which AFL (American Fuzzy Lop) is the most important and widely known \cite{zalewski2017afl}. Many state-of-the-art fuzzing tools are based on modifications and implementations of AFL \cite{fioraldi2020afl++, pham2019smart, bohme2017directed}.

Typical fuzzing tools focus only on variations of the input file and standard input. When the target program is a command-line tool, only fixed command-line options can be used for testing. Research shows that using multiple command-line options simultaneously can significantly increase test coverage; this is called "Multi-Argument Fuzzing" or "Multiple Options Fuzzing"  \cite{wu2020sq,tseng2023wei,lu2022guan,tsai2024Multiple}. These tools provide an option configuration file describing the options the target program can use and then randomly combine these options for testing.

Current multi-argument fuzzing tools suffer from two main problems. First, they require expert intervention to help write option configuration files, which is time-consuming, labor-intensive, and prone to errors. Second, the configuration files can only describe fixed, limited options, not the complete option input space, which may make it difficult to explore certain execution paths. Although existing research has used large language models (LLMs) to assist in generating parameter configuration files \cite{tsai2024Multiple}, fine-tuning LLMs still requires manpower and time. While LLMs can generate corresponding configuration files by reading the target program's user manual, they still lack the ability to fully explore the input space. This research proposes a novel multi-argument fuzzing architecture to explore the parameter input space automatically and releases an open-source prototype.
\fi